\section{Introduction}
\label{sec:introduction}
\subsection{Problem and motivation}
Recommendation systems use interaction histories to predict which items a person
will find useful. Those histories can reveal preferences and behavior that users
may not want to disclose. Federated learning changes where training occurs: user
records remain on clients, while a server coordinates the shared model. This
separation is useful, but transmitted updates can still depend strongly on an
individual's records. Differential privacy adds a formal constraint on that
dependence. In this project the protected unit is an entire user, so protecting
one rating while exposing the rest of the user's history would be insufficient.

The practical difficulty is that privacy-preserving aggregation also changes the
learning problem. A client update is clipped before contributing to an aggregate,
and random noise is added to that aggregate. Clipping can suppress informative
directions; noise can obscure weak collaborative signals. These effects interact
with client sampling, local training, the representation of items, and persistent
user embeddings. A useful experimental study must distinguish these mechanisms
instead of treating every difference from a centralized model as a privacy cost.

The project's starting hypothesis was that reducing the number of shared item
parameters could improve the privacy--utility--communication trade-off. A full
item matrix $Q\in\R^{M\times d}$ can be written approximately through factors
$A\in\R^{M\times r}$ and $B\in\R^{r\times d}$, with $r<d$. Sharing the factors
reduces the number of communicated scalars from $Md$ to $r(M+d)$. Under isotropic
Gaussian perturbation, fewer coordinates also reduce the expected squared norm
of noise in those coordinates. Neither observation establishes that recommendation
scores become less disturbed. The model uses the product $AB$, not the separate
factor coordinates, and its predictions multiply that product by a local user
vector whose scale can itself change during noisy training.

The investigation consequently developed in stages. The initial baselines
separated centralized learning, user weighting, federation, privacy, and low rank.
Their mixed outcomes motivated controlled comparisons of parameterizations and
direct measurement of effective noise. Later experiments asked whether noise
persists through local personalization and whether diagnostic trajectory
differences depend on arbitrary factor coordinates. A final candidate examined
whether non-Gaussian margin tails matter in the saved trained states. This report
presents those stages as a connected study, including the hypotheses that failed.

\subsection{Research questions}
Five questions organize the completed work.
\begin{enumerate}[label=RQ\arabic*.]
\item \textbf{Representation and utility.} At matched privacy targets and a fixed
private training horizon, when does a low-rank shared representation outperform
a full shared matrix, and what utility is lost before privacy is introduced?
\item \textbf{Parameterization and effective noise.} Does fixing one public
factor improve prediction quality once initialization, tuning opportunity,
clipping controls, and a private popularity reference are included?
\item \textbf{Persistent personalization.} After an injected shared-model
perturbation has influenced local user states, how much prediction disturbance
remains when the shared state is restored?
\item \textbf{Validity of trajectory comparisons.} How much of a paired
two-factor trajectory difference depends on rebalancing and on the rule used to
couple random noise between equivalent factor representations?
\item \textbf{Ranking-tail relevance.} In the trained states, does the exact
conditional law of a noisy item-pair margin differ enough from Gaussian
approximations to motivate a different ranking method?
\end{enumerate}
These questions concern related but distinct quantities. A difference in score
RMS does not automatically imply a difference in ranking relevance; a rank-set
change need not lower NDCG; and a smaller transmitted tensor does not establish
lower communication to convergence. Each stage therefore uses a measurement
appropriate to the claim it can support.

\subsection{Completed contributions}
The first contribution is a reproducible experimental framework: deterministic
data preparation, full-catalog ranking, explicit user-level aggregation and
accounting, matched controls, repeated training seeds, and an immutable evidence
trail. This framework supports a qualified empirical result: low rank is useful
in some strong-privacy comparisons, but it is not a general remedy for private
collaborative recommendation.

The second contribution is a controlled recommendation-specific evaluation of
effective noise. The study relates factor-space noise to the item product, local
user state, clipping, ranking, and communication. The fixed-factor method is a
baseline based on established ideas, not a newly invented algorithm. A separately
calibrated user-level private popularity method exposes the weakness of several
apparently promising comparisons against a deteriorating collaborative reference.

The third contribution is an audited diagnostic finding. Large paired distances
in rebalanced two-factor trajectories can depend strongly on coordinate coupling.
Alignment and continuation-without-balancing controls substantially narrow the
interpretation of the initial local-memory observation. This finding reproduces
on a second MovieLens dataset, although that does not establish independent user
populations or generalization to a different service.

Finally, the project documents a rejected mechanistic candidate. A correct
non-Gaussian conditional margin law produces negligible practical approximation
error in the selected real checkpoints. Retaining this result prevents a
mathematically interesting construction from being presented as an explanation
for utility loss without empirical support. Together, these contributions form
a substantial undergraduate research study. Establishing a distinct publishable
algorithm or a novel general theorem remains future work.

\subsection{Organization}
\Cref{sec:related,sec:privacy,sec:data} establish prior work, the threat model, and
the evaluation protocol. \Cref{sec:methods,sec:design} describe model variants and
experimental controls. \Cref{sec:baselines,sec:e1,sec:e2,sec:e3,sec:e4} report the
completed stages. \Cref{sec:discussion} synthesizes the evidence and limitations.
The appendices contain derivations, algorithm details, material revisions, and
reproduction instructions. The supplement retains the expanded numerical record.
